\section{The bit-to-word reduction}
\label{sec:reduction}
\begin{theorem}[Uniform bit-to-word reduction]
\label{thm:bit-word}
If one deterministic multitape machine transposes every dyadic square bit
matrix in time linear in its input length, then one such machine transposes
every canonical square matrix with entry width $b(d)$ in linear time:
\[
                              \BIT\Longrightarrow\WORD.
\]
Both problems use the raw-input model of \cref{sec:model}. The assumed bit
machine need not have an oblivious schedule or a workspace-reset guarantee.
\end{theorem}

The Lean endpoint is
\lean{BitWord.linearTimeTransposable_of_dyadicBitOracle}. The remainder of
this section outlines its proof, from address identities to native execution.

\subsection{A constant-length address factorization}
For a square of dyadic side $D=2^h$ with $h\geq1$, set
\[
 u=\ceil{\log_2 h},\quad e=u\bmod2,\quad w=u+e,\quad
 N=2h+w=2n,\quad W=2^w.
\]
The even choice of $w$ makes $n=h+w/2$ an integer and the padded
bank length a square power of two. Pad every $h$-bit entry with trailing zeros
to width $W$. The padded bank has $2^N$ bits, and each address has fields
\[
                 (\text{row}[h],\text{column}[h],\text{lane}[w]).
\]
The desired permutation exchanges the row and column fields and preserves
every lane. The bound $W<4h$ limits the padding cost.

Let $S_k$ exchange the last two $k$-bit address fields, fixing all earlier
positions; it is legal when $2k\leq N$. Let $R$ move the last address bit
to the front, and let $R^{-1}$ reverse this rotation. These are source-motion
permutations: a bit initially at address $a$ moves to the transformed address.
Composition lists below are chronological.

\begin{lemma}[At most thirteen stages]
\label{lem:thirteen}
For every $h\geq1$, a specified list of at most thirteen legal operators
$S_k,R,R^{-1}$ exchanges the two leading $h$-bit fields and fixes all tail
positions. All indices in the list are explicit functions of $h$.
\end{lemma}

\begin{proof}[Proof outline]
For $h<32$, one has $w\leq6$. Use $w$ copies of $R$, then $S_h$, then
$w$ copies of $R^{-1}$. This uses $2w+1\leq13$ stages.

For $h\geq32$, the chosen parameters satisfy $3w\leq n$. Define the
five-stage word
\[
                   Q=(S_n,S_{n-w},S_w,S_{n-w},S_n).
\]
It exchanges the first and last $w$ address positions and fixes the middle.
Conjugating the suffix exchange by this endpoint exchange gives
\[
                            (Q,S_h,Q),
\]
which has eleven stages. To check the conjugation, split the leading
fields as $A=A_0A_1$ and $B=B_0B_1$, where $|A_0|=|B_0|=w$,
and write the tail as $C$. The three operators act by
\[
\begin{aligned}
 (A_0,A_1,B_0,B_1,C)&\xrightarrow{Q}(C,A_1,B_0,B_1,A_0)\\
 &\xrightarrow{S_h}(C,B_1,A_0,A_1,B_0)\\
 &\xrightarrow{Q}(B_0,B_1,A_0,A_1,C).
\end{aligned}
\]
This is the desired head-field swap. The endpoint exchange is its own
inverse. The finite-address identities, including every
boundary case, are checked in \lean{five_suffix_end},
\lean{elevenSuffix_eq_head}, and \lean{canonicalGates_correct}. Their bank
interpretation is \lean{canonical_bank_action}.
\end{proof}

\subsection{Native execution and workspace reuse}
A suffix stage $S_k$ partitions the bank into consecutive blocks of $4^k$
bits. Each block is a raw $2^k\times2^k$ bit matrix, so one call to the bit
transposer realizes the required suffix exchange. Calls are sequential,
reusing one fixed collection of work tapes. The number of calls is not
bounded by thirteen: thirteen bounds the number of whole-bank stages.
The block lengths sum to the bank length, so linear cost per block gives
linear cost per stage.

The oracle may branch on data, move its heads adaptively, and use nonlinear
work encodings. The wrapper interprets its actual transition function on a
copied raw block and captures the append-only output on a work tape. The cleanup follows the history-and-replay method of
Bennett~\cite{Bennett1973},
with a checked finite controller for the present model. To restore the workspace,
a fixed finite record alphabet stores the old scanned
symbols and control label of each simulated microstep. Replaying these records
backward restores all unretained tapes and head positions while retaining the
captured output. Record creation, replay, erasure, and the return to the
origin are included in the bound. The alphabet and number of log tracks may
depend on the fixed oracle machine; they do not depend on the block size.

The rotation $R$ is a stable even/odd split: read the bank once, separate
the two parities, rewind, and concatenate. The inverse interleaves its two
halves. The checked split and interleave subroutines initialize and erase their
scratch tapes and restore all heads. A physical ruler controls width
expansion, cropping, and block traversal. Moving and resetting a ruler is
charged to the data it accompanies. A dimension or a loop index written in
binary is not treated as a constant-time arithmetic register.

\subsection{Uniformity and arbitrary dimensions}
The wrapper has one fixed controller. Its tape count, alphabet and control
states may depend on the chosen bit machine, but not on the input dimension.
The input contains only the matrix bits. The raw dimension loader from
\cite{SecondPaper2026} recovers $d$ and the entry width $h=b(d)$.

Set $D=2^h$ and use the lane width $W=2^{w}$ defined above. This choice
satisfies
\[
 d\leq D\leq2d,\qquad h\leq W<4h,\qquad
 V:=D^2W\leq16P,\qquad P=d^2b(d).
\]
The same procedure handles $d=1$: it pads to $D=2$, executes the
permutation, and crops. No exceptional machine or separate size advice is
needed.

The fixed controller performs the following operations in order.
\begin{enumerate}[nosep]
\item Load the raw input onto a marked work tape. Compute binary words for
$d,h,w,n$ and the gate parameters, and write thirteen physical opcode and
parameter slots. Unused slots contain an idle opcode.
\item Construct unary rulers for row counts and record widths. Expand each
entry from $h$ to $W$ bits, then pad rows and columns from $d$ to $D$.
\item Construct a half-bank ruler and execute the thirteen slots. Every
suffix slot computes its block ruler from its stored parameter and reuses
the same oracle workspace for all consecutive blocks.
\item Crop rows and columns, then crop entries to $h$ bits. Append the
result to the raw output and halt.
\end{enumerate}
The descriptor generator uses a fixed finite table for $h<32$ and binary
arithmetic for $h\geq32$. Both branches belong to the same native program.
Sharing a physical port between two subroutines is a fixed wiring choice in the controller; it neither copies
data nor changes its location during execution.

The setup trace \lean{NativePrepare.complete_trace} costs at most $40000P$.
Each padding or cropping pass costs at most $200V+200$ by
\lean{NativeNormalize.complete_trace}. The half-bank ruler has cost linear
in $V$ and its binary exponent length. Each whole-bank gate costs at most
\[
                  (\kappa+300)V+24\ell+100,
\]
where $\kappa$ depends only on the fixed bit machine and its linear-time
constant, and $\ell$ is the stored parameter-word length.
\lean{NativeStages.complete_trace} executes all thirteen slots and includes
their linking transitions. The generated parameter words have length at most
$2h+2$, so these costs sum to $O(P)$. This accounting includes physical
counter construction, repeated block calls, cleanup and raw output; it does
not charge a binary arithmetic operation as a single tape transition.

The semantic composition uses exact serialization identities. Padding commutes
with transpose, as does widening every entry by trailing zeros. The checked
address action therefore transforms the padded serialization into the padded
serialization of the transposed original matrix. Cropping recovers precisely
{\small\texttt{Transpose.encode d (Transpose.entryBits d) (Transpose.transpose A)}}.
The complete trace is \lean{NativeReductionMachine.complete_trace}.
It proves the bound $(200000+208\kappa)P$ from blank work tapes, with the exact
raw transposed output and actual halting. The compiler bridge identifies its
initial configuration with the original machine model's initial state and
yields \lean{NativeReductionMachine.native_transpose}, from which
\cref{thm:bit-word} follows.
